\documentclass[11pt]{article}

\usepackage[margin=1in]{geometry}
\usepackage{amsmath,amssymb,amsthm,mathtools}
\usepackage{stmaryrd}
\usepackage[expansion=false]{microtype}
\usepackage[hidelinks]{hyperref}

%% ---- theorem environments --------------------------------------------
\theoremstyle{plain}
\newtheorem{theorem}{Theorem}
\newtheorem{proposition}{Proposition}
\newtheorem{lemma}{Lemma}
\newtheorem{corollary}{Corollary}
\theoremstyle{definition}
\newtheorem{definition}{Definition}
\newtheorem{remark}{Remark}
\newtheorem{obligation}{Obligation}

%% ---- macros ----------------------------------------------------------
\newcommand{\den}[1]{\llbracket #1 \rrbracket}
\newcommand{\Name}{\mathrm{Name}}
\newcommand{\Proc}{\mathrm{Proc}}
\newcommand{\Mfin}{M_{\mathrm{fin}}}
\newcommand{\Pfin}{\mathcal{P}_{\mathrm{fin}}}
\newcommand{\rfor}[2]{\mathsf{for}(#1\!\leftarrow\!#2)\,}
\newcommand{\rout}[2]{#1\,!(#2)}
\newcommand{\comp}{\mathrel{\asymp}}            % comparable in the prefix order
\newcommand{\loc}{\mathrel{\rhd}}               % "located at": w |> Q
\newcommand{\quo}{@}
\newcommand{\der}{{\ast}}
\newcommand{\rr}{\rightharpoonup}
\newcommand{\zip}{\mathbin{\bowtie}}
\newcommand{\cut}{\chi}
\newcommand{\React}{\rho}
\newcommand{\Sig}{\Sigma}
\newcommand{\eps}{\varepsilon}
\newcommand{\bisr}{\sim_{r}}
\newcommand{\bis}{\sim}
\DeclareMathOperator{\supp}{supp}

\title{\textbf{Paths are Subspaces}\\[2pt]
\large a cut-triggered polymorphism of RSpace over path-keys,\\ and its distributive law}
\author{Lucius Gregory Meredith\\
\small CEO, F1R3FLY.io, and Chief Science Officer, F1R3FLY Industries\\
\small\texttt{lgreg.meredith@gmail.com}}
\date{}

\begin{document}
\maketitle

\begin{abstract}
The RSpace denotation of the rho calculus in the knotted universe factors into two independent
pieces: a reflection structure on the \emph{values} (the colour-swap $s$, with $\quo=\kappa\circ s$
and $\der=s\circ\kappa^{-1}$) and a store monoid on the \emph{keys}
($\Name\rr\Mfin(-)$ under key-wise $\uplus$). Anything done to the key object therefore leaves the
reflection law untouched, so RSpace is polymorphic in its key type at no cost on the reflection side.
We take the keys to be \emph{paths}, $K=\Name^{\ast}$, the free monoid on names, whose store is a
trie. The prefix order on paths supplies a notion of one channel lying in the \emph{subspace} of
another, and we read communication accordingly: a consume at $\mathrm{path}_1$ and a produce at
$\mathrm{path}_2$ react when the paths are comparable, not merely equal. We insist throughout on a
single discipline --- \emph{no implicit interaction}: the store never reacts with itself; a reaction
is triggered only by a cut, an explicit $\mid$ co-locating a consume and a produce, exactly as in
the source calculus. What comparability changes is not \emph{when} a redex exists but \emph{what} a
given cut matches and \emph{what residual} it leaves; and because the residual of a deep cut is again
a single cut, one interaction can cause nested interactions. We give COMM in this generality as a
pair of dual rules, exhibit the cascade as the reflection law $s\circ s=\mathrm{id}$ unsealing a
located sub-RSpace, and write the dynamics as an explicit distributive law of the cut over the store
monad, a bialgebra in the Turi--Plotkin sense, in which the unit axiom \emph{is} the
no-implicit-interaction principle. We then check the ledger: congruence by construction, the exact
monoid homomorphism, OSLF adequacy and full abstraction, and the source paper's risk accounting all
survive, with comparability the only new ingredient in the zip. Where the paths share prefixes ---
which, because names are quoted processes, is whenever channels share subterm structure --- the trie
is a radix-compressed, structurally shared, Merkle-addressable store, and we record where each
compression pays off.
\end{abstract}

\section{Introduction}

The companion note \cite{KnotRho} models the reflective higher-order (rho) calculus
\cite{MR05} inside the knotted universe of \cite{Knot}, with the quote $\quo$ reading as the
colour-swap that seals a process into an atom of the opposite colour and the dereference $\der$ as
unsealing it. The dynamics there factor through an abstraction, \emph{RSpace}: parallel composition
is interpreted as keyed multiset union, so the denotation $\den{-}$ is a monoid homomorphism and the
parallel laws of structural congruence become literal equalities of tables, hashtables from channels
to polarity-homogeneous bags of produces and consumes; communication is eager reaction, the zip of
dual bags at a key, and RSpace is the store monad together with this zip, a bialgebra --- the precise
content of ``strictly more than a monad''.

The keys of that store are names, and in rho a name is a quoted process; the keys therefore have
structure. This note studies one polymorphic interpretation of RSpace obtained by varying that
structure. We make the key object a \emph{path}: $K=\Name^{\ast}$, the free monoid on names, each key
a route through a trie to a sub-RSpace. The interaction algebra of RSpace acts on entire RSpaces, so
once keys carry a path order one may ask reaction to fire not only at equal keys but at
\emph{comparable} ones --- a consume in the subspace addressed by a prefix of the produce's path, or
\emph{vice versa}.

One design decision governs the development.

\paragraph{No implicit interaction.}
The store is inert. It never reacts with itself, and there is no ambient broadcast in which a stored
produce finds every comparable consume in the trie. A reaction is triggered only by a \emph{cut}: an
explicit parallel composition $\mid$ that co-locates a consume and a produce, exactly as $K=\mid$ is
the single base constructor heading the left-hand side of COMM in the source calculus. What the path
order buys is not new triggers but a richer \emph{matching} at an existing trigger:
$\rfor{y}{\mathrm{path}_1}P \mid \rout{\mathrm{path}_2}{Q}$ is a local cut even when
$\mathrm{path}_1$ and $\mathrm{path}_2$ are not equal on the nose, and the comparability of the two
paths decides what is delivered and what residual remains. Because the residual of a deep cut is
again a single cut, an interaction can \emph{cascade}: it can manufacture nested interactions, each
of which is itself a cut, with nothing implicit in between.

This decision is what keeps the construction conservative. We show below that it leaves the locality
of the cut --- the reactive contexts are still $[\,\cdot\,]\mid R$ --- and hence the
congruence-by-construction argument, the exact homomorphism, OSLF adequacy, and the risk ledger of
\cite{KnotRho} all intact; the only genuinely new object is the comparability-aware zip, which we
write out as a distributive law of the cut over the store monad. We assume throughout the modern
presentation of \cite{KnotRho}: $\quo P$ for quoting, $\der x$ for dereference, $\rfor{y}{x}P$ for
input, $\rout{x}{Q}$ for output, $\equiv$ for structural congruence, $\bis$ for context bisimilarity,
$\bisr$ for resource bisimilarity, and $\den{-}\colon\Proc\to T(\Proc)$ for the RSpace denotation
into tables $T(R)=\Name\rr\bigl(\Mfin(R)\uplus\Mfin([N]R)\bigr)$.

\section{RSpace, parametric in the key object}\label{sec:param}

\subsection{The factoring lever}

Write RSpace's table functor with its key object exposed as a parameter:
\[
  T_{K}(R)\;=\;K\;\rr\;\bigl(\Mfin(R)\;\uplus\;\Mfin([N]R)\bigr),
\]
a finite partial map from keys to polarity-homogeneous bags of payloads --- a key holds either a
multiset of output payloads (produces, data) or a multiset of input continuations (consumes,
abstractions $(y)P$), never a mixture. The source paper is the instance $K=\Name\cong\Proc$, the
keys quoted processes. The construction splits cleanly:

\begin{remark}[The reflection/store factoring]\label{rem:factor}
RSpace is two independent structures glued at the table. The first is a reflection on the
\emph{values}: the colour-swap $s\colon\mathcal S\cong\overline{\mathcal S}$, $s\circ s=\mathrm{id}$,
with $\quo=\kappa\circ s$ and $\der=s\circ\kappa^{-1}$, validating $\der\quo P\equiv P$ as the
involutivity of $s$. The second is a store monoid on the \emph{keys}: $K\rr\Mfin(-)$ with
multiplication $\uplus$, the free commutative monoid lifted key-wise. The reflection law concerns
names, not the store monoid, and is undisturbed by any choice of $K$; the store laws concern keys,
not values, and are undisturbed by any choice of value-reflection. Consequently every key-object
polymorphism of RSpace validates Proposition (reflection) and clause~(i) of the adequacy theorem of
\cite{KnotRho} \emph{verbatim}, for free. The entire content of changing $K$ lives in the
store-and-zip factor.
\end{remark}

\subsection{What a key object must supply}

The minimal interface to instantiate $T_{K}$ as a store is decidable equality on $K$: enough to
detect that a key carries both polarities, i.e.\ to recognise a redex. To do \emph{path} reaction one
asks for more --- the free monoid $(\Name^{\ast},\cdot,\eps)$ and its prefix order. To retain the
congruence-by-construction of \cite{KnotRho} one asks, finally, that the category of $K$-contexts
possess the relative pushouts of Leifer--Milner \cite{LM00,Sewell02}, adapted to GSLTs \cite{GSLT};
\S\ref{sec:transfer} shows the no-implicit-interaction discipline is exactly what makes this
no harder for paths than for bare names.

\section{Paths and the trie}\label{sec:trie}

\subsection{The path key object}

Fix $K=\Name^{\ast}$, finite sequences of names, with $\eps$ the empty path, $\cdot$ concatenation,
$p\sqsubseteq q$ the prefix order ($q=p\cdot w$ for some suffix $w$), and $p\wedge q$ the
longest common prefix, which exists for all $p,q$ because the prefix order is a tree order. Write
\[
  p\comp q \;:\Longleftrightarrow\; p\sqsubseteq q \ \text{ or }\ q\sqsubseteq p
\]
for \emph{comparability}: one path lies on the route to the other. (We reserve $\#$ for the freshness
of \cite{KnotRho}.) The monoid acts on tables by descent: for $p\in K$ and a table $u$, the
\emph{graft} $p\cdot u$ is the table with $(p\cdot u)(p\cdot k)=u(k)$ and undefined off
$p\cdot K$, and the \emph{residual} $u/p$ is the sub-table rooted at $p$, $(u/p)(k)=u(p\cdot k)$. A
located output is sugar: $w\loc Q$ is the process with
\[
  \den{w\loc Q}\;=\;\bigl[\,w\mapsto\{\!|\,\den Q\,|\!\}^{+}\,\bigr],
  \qquad \eps\loc Q\equiv Q,
\]
i.e.\ $Q$ planted at sub-path $w$; since names are quoted processes, $\quo(w\loc Q)$ is an honest
name. Dually $\eps$ is a transparent rendezvous: an input or output on the empty path is on the node
itself.

\subsection{The store is a trie, and the trie compresses}

Modulo $\equiv$ the parallel fragment is the free commutative monoid, and interpreting $\mid$ as
key-wise $\uplus$ over the index $\Name^{\ast}$ keeps $\den{-}$ a monoid homomorphism on the nose
exactly as before. The store $\Name^{\ast}\rr\Mfin(-)$ is, with prefix sharing, a trie:
\[
  \mathrm{Trie}(V)\;\cong\;\bigl(\Name^{\ast}\rr V\bigr),
  \qquad
  \mathrm{Trie}(V)\;\cong\;V\times\bigl(\Name\rr\mathrm{Trie}(V)\bigr),
\]
the isomorphism on the left being precisely the statement that prefix sharing is sound: a trie is the
representation of a path-indexed map in which common prefixes are stored once. The recursive
presentation on the right makes every internal node itself a full RSpace --- its own bag plus its
children --- which is what licenses ``interaction acts on entire RSpaces'' for nested stores.

\begin{remark}[Four readings]
The path structure admits, in increasing order of how much of the monoid the reaction sees:
\emph{(0)} paths as opaque compound names --- reaction at equal full paths, the trie a storage
representation only, every theorem of \cite{KnotRho} transferring verbatim; \emph{(1)} the nested
trie store with exact-key reaction --- same behavioural and resource equivalences, but the recursion
now exposed in the type and $\uplus$ realised as recursive trie-merge; \emph{(2)} subspace reaction
--- a cut whose paths are comparable matches across the prefix order; and \emph{(3)} meet reaction
--- a cut rendezvous at $p\wedge q$ even for incomparable siblings. Readings 0 and 1 are
representational. Readings 2 and 3 change the matching, and are the subject of
\S\ref{sec:subspace}--\S\ref{sec:dist}. We develop the subspace reading in full and treat the meet
reading as its symmetric extension.
\end{remark}

\section{Cut-triggered subspace reaction}\label{sec:subspace}

\subsection{COMM in the subspace, as a pair of dual rules}

Communication is triggered by a cut and nothing else. When the two paths of a cut are comparable, the
single COMM rule of the source calculus,
\[
  \frac{x\equiv_{N}z}{\rfor{y}{x}P \mid \rout{z}{Q} \longrightarrow P\{\quo Q/y\}}\;\textsc{Comm},
\]
splits according to which side is the deeper. Let the input be on $\mathrm{path}_1$ and the output on
$\mathrm{path}_2$.

\begin{definition}[Subspace COMM]\label{def:subcomm}
There are two dual reactions, each a local cut.
\begin{itemize}
\item[\textsc{(Out-deep)}] If $\mathrm{path}_1\sqsubseteq\mathrm{path}_2$, write
  $\mathrm{path}_2=\mathrm{path}_1\cdot w$. The output lies in the subspace named by the input; the
  input receives the still-located sub-RSpace, sealed:
  \[
    \rfor{y}{\mathrm{path}_1}P \mid \rout{\mathrm{path}_2}{Q}
    \;\longrightarrow\;
    P\{\,\quo(w\loc Q)\,/\,y\,\}.
  \]
\item[\textsc{(In-deep)}] If $\mathrm{path}_2\sqsubseteq\mathrm{path}_1$, write
  $\mathrm{path}_1=\mathrm{path}_2\cdot w$. The output arrives at the shallower channel and the input,
  waiting deeper, must descend $w$ into the payload; the cut \emph{unpacks} into a fresh, deeper cut:
  \[
    \rfor{y}{\mathrm{path}_1}P \mid \rout{\mathrm{path}_2}{Q}
    \;\longrightarrow\;
    \rfor{y}{w}P \mid Q.
  \]
\end{itemize}
Both collapse to \textsc{Comm} at $w=\eps$: \textsc{(Out-deep)} gives $P\{\quo Q/y\}$ since
$\eps\loc Q\equiv Q$, and \textsc{(In-deep)} gives $\rfor{y}{\eps}P\mid Q\longrightarrow P\{\quo Q/y\}$
by the transparency of the root rendezvous. The propagation rules \textsc{Par}, \textsc{Struct} are
unchanged.
\end{definition}

The right-hand side of \textsc{(In-deep)} is again a single parallel composition, a local cut, not a
probe of the store. This is the precise sense of ``an interaction can cause nested interactions'':
one cut produces a residual that is itself a redex, fired by the same rules, with nothing ambient in
between. Whether it cascades depends on $Q$: if $Q$ outputs on a channel comparable to $w$ the new
cut fires; otherwise it is a pending input, correctly stuck.

\subsection{Reflection is the cascade engine}

The seal of \textsc{(Out-deep)} is where reflection does double duty.

\begin{proposition}[Cascade via $s\circ s=\mathrm{id}$]\label{prop:cascade}
In \textsc{(Out-deep)} the receiver binds the name $\quo(w\loc Q)=\kappa\,s(w\loc Q)$, the located
sub-RSpace sealed by the swap. Wherever the continuation $P$ dereferences it, $\der$ unseals it,
\[
  \der\,\quo(w\loc Q)\;=\;s\,\kappa^{-1}\kappa\,s\,(w\loc Q)\;=\;s\,s\,(w\loc Q)\;=\;w\loc Q,
\]
splicing the located output back into the body of $P$, where it may cut with $P$'s own consumes one
level deeper. The colour-swap therefore fires twice per cascade step: once as the seal
$\quo=\kappa\circ s$ packaging the residual sub-RSpace into the delivered name, once under
$\der=s\circ\kappa^{-1}$ unsealing it for the next reaction.
\end{proposition}

The cascade is thus not a new nonlocal mechanism but the reflection-through-reflection recursion the
source calculus already relies on for recursion without a primitive replication operator (the server
$D\equiv\rfor{y}{x}(\der y\mid\rout{x}{y})\mid\rout{x}{\quo(\cdots)}$ of \cite{KnotRho}), now also
threading the trie.

\section{The distributive law}\label{sec:dist}

We now present the dynamics of Definition~\ref{def:subcomm} as a distributive law of the cut over the
store monad, in the bialgebra format of Turi--Plotkin \cite{TP97}. Two facts orient the construction:
the store monad is untouched (still $\Name^{\ast}\rr\Mfin(-)$ with $\uplus$), and the no-implicit-interaction
principle is, formally, the \emph{unit} axiom of the law.

\subsection{The store monad and the pre-tables}

Let $(T,\eta,\mu)$ be the store monad, $\mu=\uplus$, $\eta$ the singleton placement, on the value
functor of \S\ref{sec:param}. During reaction a key may transiently carry both polarities, so we work
on \emph{pre-tables}
\[
  \overline T(R)\;=\;K\;\rr\;W(R),
  \qquad
  W(R)\;=\;\Mfin(R)\times\Mfin([N]R),
\]
a key holding a produce-bag $B^{+}$ and a consume-bag $B^{-}$; a table is the polarity-homogeneous
sub-object (one bag empty at every key). The cut installs both by key-wise union: $\den{P\mid Q}=
\den P\uplus\den Q$, with $\uplus$ pure (no firing) so that $\den{-}$ stays a function and reaction is
the separate dynamics.

\subsection{The fire of one matched pair}

A \emph{consume item} is $(p,(y)P)$ and a \emph{produce item} is $(q,Q)$; they form a redex iff
$p\comp q$. Firing a matched pair returns the denotation of the corresponding rule's right-hand side,
to be merged back by $\uplus$.

\begin{definition}[Fire]\label{def:fire}
For $p\comp q$ set $w$ by $q=p\cdot w$ (when $p\sqsubseteq q$) or $p=q\cdot w$ (when $q\sqsubseteq p$),
and define
\[
  \mathrm{fire}\bigl((p,(y)P),(q,Q)\bigr)\;=\;
  \begin{cases}
    \den{P\{\,\quo(w\loc Q)/y\,\}} & p\sqsubseteq q \quad(\textsc{Out-deep}),\\[4pt]
    \den{\rfor{y}{w}P}\;\uplus\;\den{Q} & q\sqsubseteq p \quad(\textsc{In-deep}).
  \end{cases}
\]
At $p=q$ both lines equal $\den{P\{\quo Q/y\}}$, the source paper's COMM. The \textsc{In-deep} line is
a fresh cut: $\den{\rfor{y}{w}P}=[\,w\mapsto\{\!|(y)\den P|\!\}^{-}]$ unioned with $\den Q$, eager
reaction continuing within.
\end{definition}

\subsection{The cut term and the law}

Reaction enters only through juxtaposition. Given pre-tables $s,t$, the cut formed by composing them
is the eager one-pass cross-reaction of all comparable opposite-polarity pairs with one side from
each operand.

\begin{definition}[Cut term]\label{def:cutterm}
\[
  \cut(s,t)\;=\;
  \biguplus
  \Bigl\{\,
    \mathrm{fire}(c,d)\;:\;
    c\ \text{a consume item of }s,\ d\ \text{a produce item of }t,\ c\comp d
  \Bigr\}
  \ \uplus\ (s\leftrightarrow t),
\]
where $(s\leftrightarrow t)$ is the symmetric term (consumes of $t$ against produces of $s$), and at
each key the matched pairs are zipped along chosen list orderings, the surplus of the longer bag
remaining --- the located analogue of the zip of \cite{KnotRho}, its pairing nondeterminism the
source calculus's choice of which send meets which receive, now ranging over comparable rather than
equal keys.
\end{definition}

Let $\React\colon\overline T\to\overline T$ denote eager reaction to a (subspace-)polarity-homogeneous
normal form: fire every comparable opposite-polarity pair, including those created by merged products,
until none remains. The bialgebra content is two equations.

\begin{definition}[Distributive law of the cut over the store]\label{def:dist}
Let $\Sig$ be the signature functor whose binary component is the cut $\mid$. A distributive law
\[
  \lambda\;\colon\;\Sig\,T\;\Longrightarrow\;T\,\Sig
\]
of $\Sig$ over the monad $(T,\eta,\mu)$ is a natural transformation satisfying
\[
  \text{\emph{(D1, unit)}}\quad \lambda\circ\Sig\eta=\eta\Sig,
  \qquad
  \text{\emph{(D2, mult.)}}\quad \lambda\circ\Sig\mu=\mu\Sig\circ T\lambda\circ\lambda T.
\]
On the cut component $\lambda$ is realised by $\uplus$-with-cross-reaction: on a pair of stores
$(s,t)$ it returns their key-wise union enriched by the cut term, $s\uplus t\uplus\cut(s,t)$, with the
located residuals of Definition~\ref{def:fire} reinstalled and reaction iterated to the normal form.
\end{definition}

\begin{proposition}[The unit axiom is no-implicit-interaction]\label{prop:unit}
$\cut(s,\varnothing)=\varnothing=\cut(\varnothing,t)$, hence $\lambda$ on a cut against the empty
store is the bare injection: $\React(s\uplus\varnothing)=\React(s)$. Equivalently \emph{(D1)}: cutting
against a unit store causes no reaction. A reaction requires a cut that actually co-locates opposite
polarities; the store never reacts with itself.
\end{proposition}

\begin{proof}[Proof sketch]
$\cut(s,t)$ ranges over pairs $(c,d)$ with $c$ from one operand and $d$ from the other; with one
operand empty there are no such pairs, so the cut term is empty and $\lambda$ reduces to
$\uplus$, whose unit is $\varnothing$. This is precisely the statement that $\Sig\eta$ followed by
$\lambda$ equals $\eta\Sig$: composing a process with the inert unit and then ``reacting'' is the same
as not reacting.
\end{proof}

\begin{proposition}[Multiplication, up to resource bisimilarity]\label{prop:mult}
For pre-tables $s,t$,
\[
  \React(s\uplus t)\;\bisr\;\React\bigl(\React(s)\uplus\React(t)\uplus\cut(s,t)\bigr).
\]
The right-hand side is the \emph{(D2)} pattern: react the operands, merge, and add exactly the
cross-cuts the union created; the cross-cuts are the only reactions not already internal to $s$ or
$t$, and they are exactly the comparable opposite-polarity pairs straddling the cut. Iterating
$\React$ closes the cascade.
\end{proposition}

\begin{remark}[What kind of object this is]
The law does not change the monad: $\uplus$ is verbatim the store multiplication, and
Proposition~13 of \cite{KnotRho} (the exact homomorphism, the parallel laws as literal table
equalities) holds unchanged because $\uplus$ never fires. RSpace remains ``store monad $+$ zip''; the
zip is merely comparability-aware. The reflection law is likewise untouched, by
Remark~\ref{rem:factor}. The cascade is the inductive closure of $\lambda$ over its own
\textsc{In-deep} residuals; it terminates whenever the source process is finite, because each fire
strictly consumes a matched pair, and the only route to nontermination is reflection looping --- the
single $\omega$ of \S\ref{sec:transfer}, not a new infinity.
\end{remark}

\subsection{The meet variant}

The pure-meet reading fires at $p\wedge q$ even for incomparable siblings, $p=m\cdot u$ and
$q=m\cdot v$ with $u,v$ both nonempty, $m=p\wedge q$. It remains cut-triggered. It carries one genuine
degree of freedom the subspace reading does not: with neither path containing the other one must fix a
convention for both residual suffixes, the natural symmetric move delivering each side's residual into
the other and relocating the continuation under $m$. This is where the richest cascades live, and
where one owes an explicit relocation operator; the subspace reading keeps a clean
container/contained relationship and is the conservative one to develop first.

\section{What transfers}\label{sec:transfer}

We record, against the apparatus of \cite{KnotRho}, what the subspace dynamics preserves. The thread
throughout is locality: because reaction is cut-triggered, the reactive contexts do not change.

\paragraph{Congruence, by construction.}
The reactive contexts are still $[\,\cdot\,]\mid R$ --- a cut always lies at a parallel composition,
including the manufactured right-hand cut of \textsc{(In-deep)} --- so the category of contexts and
its idem-pushouts never see the matching predicate. The relative-pushout hypothesis (Leifer--Milner
\cite{LM00,Sewell02}, GSLT minimal-context \cite{GSLT}) is no harder than for bare names, and context
bisimilarity is a congruence by construction: $P\bis Q\Rightarrow F[P]\bis F[Q]$. The
action-labelled closure-under-substitution question does not arise.

\paragraph{The label set, three views.}
The labels are still the parallel cuts $\partial T_{K}$; the OSLF-generated Hennessy--Milner
modalities $\langle F\rangle\varphi$ are still indexed by the families
$[\,\cdot\,]\mid\rout{\mathrm{path}_2}{Q}$, now with $\mathrm{path}_2$ ranging over the comparables of
the input path --- a larger, but still context-shaped, family. Firing, Witnessing and Placing remain
three readings of one label set, and OSLF adequacy gives
$P\bis Q\Leftrightarrow\forall\varphi.(P\models\varphi\Leftrightarrow Q\models\varphi)$.

\paragraph{The homomorphism and the bialgebra shape.}
$\uplus$ is untouched, so the exact monoid homomorphism survives and the parallel laws are literal
table equalities; reaction is a richer zip over the same store monad, the ``store monad $+$ zip''
bialgebra of \cite{KnotRho} with comparability-matching and located residuals the only changes. The
polarity-homogeneity invariant of reduced tables generalises to its subspace form --- no two
\emph{comparable} keys carry opposite polarities --- and no-implicit-interaction is exactly the claim
that eager reaction maintains this invariant: every opposite-polarity comparable pair was introduced
by a $\mid$, so firing it is sound and traceable to a cut.

\paragraph{Full abstraction and the risk ledger.}
The factoring of \cite{KnotRho} onto $\nu B$ persists: resource bisimilarity (counting located
copies) refines context bisimilarity, with the projection $\pi$ forgetting multiplicity and keys.
Full abstraction $\den{P}_{B}=\den{Q}_{B}\Leftrightarrow P\bis Q$ holds for a congruence with no
closure step. And the risk ledger is undisturbed: each cut is finite, a cascade from finite terms is
finite, and the only source of nontermination is reflection looping, the same single $\omega$ ---
black infinity, an $\omega$-fold multiplicity --- the source paper already charges. The path machinery
spends no new infinity; it routes the existing reflective recursion through the trie.

\section{Where the trie compresses}\label{sec:compress}

Four mechanisms, in increasing order of how much they touch the semantics rather than the
representation.

\emph{Radix (Patricia) collapse} \cite{Morrison68}. Stores are sparse, so the trie is mostly chains
of single-child nodes; collapsing each chain into one edge labelled by a path segment takes the depth
from one node per name to one node per branch point. A storage win, available in every reading.

\emph{Structural sub-trie sharing (hash-consing, copy-on-write).} Eager normal forms are
subspace-polarity-homogeneous: a reduced sub-trie is inert, with no pending redex, hence shareable by
reference and duplicated lazily only when one side reacts. This is where the recursive trie-merge of
Reading~1 pays, and it is the operational face of the resource layer: copy-on-write is ``do not
materialise a multiplicity until it is observed.''

\emph{Merkle addressing.} Tag each node with the hash of its children and sub-RSpaces become
content-addressed. Resource bisimilarity (a key-by-key, multiplicity-matched comparison of tables)
then short-circuits equal subtrees in $O(1)$ --- equal root hash means $\bisr$-equal sub-store, no
descent --- and content addressing is native, not bolted on, because $\Name=\quo\Proc$: the key
already \emph{is} a hash of a quoted process.

\emph{The path is the structure of the name.} Since names are quoted processes, the natural path is
not an arbitrary string but the spine of the quoted process itself --- its parallel decomposition,
its descent through the term. Then content-addressed channels that share \emph{subterms} share
\emph{trie structure}, automatically, by the structure of the quote; and large quoted-process names,
the case where opaque keys waste the most, are exactly where the de-duplication is greatest. This is
the denotational reading of what the on-chain RSpace history store already computes when it persists
to a Merkle radix trie keyed by hashes of channels: Readings~1--3 are the semantics that
implementation has been evaluating all along.

\section{Obligations}

In the manner of \cite{KnotRho}, \S 8, the points a fully rigorous development would discharge.

\begin{obligation}[The distributive law]
Verify \emph{(D1)} and \emph{(D2)} for $\lambda$ of Definition~\ref{def:dist} in the Turi--Plotkin
format \cite{TP97}, with $\den{-}$ the induced bialgebra morphism. \emph{(D1)} is
Proposition~\ref{prop:unit}; \emph{(D2)} is the coherence of cut formation across nested unions, of
which Proposition~\ref{prop:mult} states the $\bisr$-form. Confirm that the located residuals of
Definition~\ref{def:fire} make $\lambda$ natural in $R$.
\end{obligation}

\begin{obligation}[Eager confluence up to $\bisr$]
A shallow consume may be comparable to several deeper produces (fan-out) and a deep input unpacks into
a payload (descent), so eager reaction must be shown confluent up to resource bisimilarity with the
new residuals: all maximal reaction sequences reach $\bisr$-equal normal forms, independent of the
pairing nondeterminism of Definition~\ref{def:cutterm} and the fan-out order.
\end{obligation}

\begin{obligation}[Substitution and relocation]
Prove $\den{P\{\quo(w\loc M)/y\}}=\den{P}\{\quo(w\loc\den M)/y\}$, the located generalisation of the
substitution lemma of \cite{KnotRho}, carried by the abstraction functor's evaluation
$[N]R\times\Name\to R$, with the seal $\kappa\circ s$ applied to the message. For the meet variant,
establish the relocation operator that delivers each incomparable residual into the other and pin down
whether it is primitive or derivable.
\end{obligation}

\begin{obligation}[Relative pushouts for path-contexts]
Discharge the locality claim of \S\ref{sec:transfer}: the category of rho contexts over path-keys,
with plugging as composition, possesses relative pushouts, so the idem-pushout labels are well defined
and the Leifer--Milner congruence theorem applies. The cut lying always at a parallel composition
(including the \textsc{In-deep} residual) and names being quoted structure rather than $\nu$-bound are
the features that should make the relative pushouts exist and be computable.
\end{obligation}

\begin{obligation}[Subspace-homogeneous normal forms]
Show eager reaction maintains the subspace-polarity-homogeneity invariant --- no two comparable keys
carry opposite polarities in a normal form --- and that this is equivalent to no-implicit-interaction:
every comparable opposite-polarity pair is traceable to a cut.
\end{obligation}

\begin{obligation}[The trie store in the knotted universe]
Encode the trie inside the knotted universe so that recursion is accommodated and the swap $s$ and
knots $\kappa,\bar\kappa$ respect the path structure and the count arithmetic, multiplicity-$\omega$
being exactly the black infinity of \S\ref{sec:transfer}; confirm $T_{\Name^{\ast}}$ preserves weak
pullbacks (the index-set change to $\Name^{\ast}$ leaves the refinement-monoid argument of
\cite{GS01} intact, the $\mathbb N\cup\{\omega\}$ completion to be re-checked).
\end{obligation}

\section{Conclusion}

Reading RSpace's keys as paths makes the store a trie and the prefix order a notion of one channel
lying in another's subspace. Insisting that reaction stay cut-triggered --- the store never reacts
with itself, an interaction requires a $\mid$, and an interaction may manufacture nested interactions
--- turns the subspace reading into a conservative extension: COMM splits into the dual
\textsc{Out-deep}/\textsc{In-deep} rules, the cascade is the reflection law $s\circ s=\mathrm{id}$
unsealing a located sub-RSpace, and the dynamics is a distributive law of the cut over the unchanged
store monad whose unit axiom is precisely the no-implicit-interaction principle. The cost is paid on
the confluence side, where it is the same genus of rewrite obligation the eager semantics already
carries; the congruence side, the exact homomorphism, OSLF adequacy and full abstraction, and the
risk ledger are all preserved. And because names are quoted processes, the path can be the structure
of the name, so the trie is not only a storage compression but the denotational content of the Merkle
radix store the implementation already runs.

\paragraph{Future work.} The meet variant and its relocation operator; the lazy reaction strategy,
treating a comparable opposite-polarity pair as a standing redex rather than reducing eagerly; and the
join-pattern axis --- paths grade the store by the free monoid $\Name^{\ast}$ under concatenation,
whereas join patterns grade it by tupling $\Name^{n}$, and the two gradings compose.

\begin{thebibliography}{99}

\bibitem{KnotRho} L. G. Meredith. \emph{Quoting is Colour-Swap: a model of the rho calculus in the
knotted universe.} Manuscript, 2026.

\bibitem{Knot} L. G. Meredith. \emph{A Knotted Universe: a new notion of reflective set theories.}
Manuscript, 2026.

\bibitem{MR05} L. Gregory Meredith and Matthias Radestock. A reflective higher-order calculus.
\emph{Electr. Notes Theor. Comput. Sci.}, 141(5):49--67, 2005.

\bibitem{GSLT} L. G. Meredith, Christian Wells, et al. \emph{Minimal-context transition systems and
the observer-specified logic functor for graph-structured lambda theories.} Manuscript, forthcoming.

\bibitem{LM00} James J. Leifer and Robin Milner. Deriving bisimulation congruences for reactive
systems. In \emph{CONCUR 2000}, LNCS 1877, pages 243--258. Springer, 2000.

\bibitem{Sewell02} Peter Sewell. From rewrite rules to bisimulation congruences. \emph{Theoret.
Comput. Sci.}, 274(1--2):183--230, 2002.

\bibitem{TP97} Daniele Turi and Gordon Plotkin. Towards a mathematical operational semantics. In
\emph{12th Annual IEEE Symposium on Logic in Computer Science (LICS 1997)}, pages 280--291. IEEE,
1997.

\bibitem{GS01} H. Peter Gumm and Tobias Schr\"oder. Monoid-labeled transition systems. \emph{Electron.
Notes Theor. Comput. Sci.}, 44(1):185--204, 2001.

\bibitem{Morrison68} Donald R. Morrison. PATRICIA --- Practical Algorithm To Retrieve Information
Coded In Alphanumeric. \emph{J. ACM}, 15(4):514--534, 1968.

\bibitem{Milner93} Robin Milner. The Polyadic $\pi$-Calculus: A Tutorial. In \emph{Logic and Algebra
of Specification}, NATO ASI Series F 94, pages 203--246. Springer, 1993.

\bibitem{Rutten00} J. J. M. M. Rutten. Universal coalgebra: a theory of systems. \emph{Theoret.
Comput. Sci.}, 249(1):3--80, 2000.

\end{thebibliography}

\end{document}
